\documentclass[../../../include/open-logic-section]{subfiles}

\begin{document}

\olfileid{his}{set}{pathology}
\olsection{Patologi}

Nanging, dhéfinisi \emph{limit} jebulé ngidinaké
sawetara konstruksi kang rada ``patologis''.

Watara dasawarsa 1830-an, Bolzano nemokaké
fungsi kang \emph{kontinu ing saben titik},
nanging \emph{ora bisa didiferensiasi ing titik endi waé}.
(Sayangé, Bolzano ora tau nerbitaké asil iki;
para matématikawan nemoni gagasan iki ing taun
1872 amarga Weierstrass uga nemokaké gagasan
kang padha kanthi mandiri.)\footnote{Sajarah
iki dicathet kanthi rinci banget ing cathetan
sikil artikel Wikipedia ngenani
\href{http://en.wikipedia.org/wiki/Weierstrass_function}{fungsi
Weierstrass}.} Bab iki, paling ora, nggumunaké.
Gampang nemokaké fungsi kaya $|x|$ kang
kontinu ing saben titik nanging ora bisa
didiferensiasi ing sawijining titik.
Nanging fungsi kang kontinu ing saben
titik lan \emph{ora ing titik endi waé}
bisa didiferensiasi iku bab kang béda
banget. Coba pikirna sedhéla carané
nggambar fungsi kaya mangkono. Supaya
kontinu, fungsi iku kudu bisa digambar
tanpa tau ngangkat pulpen saka kertas;
nanging supaya ora bisa didiferensiasi
ing titik endi waé, arah pulpen kudu
malih kanthi dadakan tanpa kendhat.

``Patologi'' liya banjur muncul. Tanggal
5 Januari 1874, Cantor nulis layang
marang Dedekind lan ngajokaké masalah iki:
\begin{quote}
Apa sawijining bidang (upamané kothak
kalebu pinggiré) bisa dipasangaké
siji-siji karo sawijining garis
(upamané garis lurus kalebu
rong pucuké), saéngga saben titik
ing bidang cocog karo siji titik
ing garis, lan kosok baliné saben
titik ing garis cocog karo siji
titik ing bidang?

Nganti saiki wangsulané pitakonan
iki isih katon angel banget
kanggo aku---senajan ing kéné
uga wong kaya-kaya kepéngin
mangsuli \emph{ora} nganti
buktiné katon meh ora perlu.
[Dikutip ing \citealt{Gouvea2011}]
\end{quote}
Nanging ing taun 1877, Cantor mbuktèkaké
yèn pangirané salah. Satemené, garis
lan kothak nduwèni cacah titik kang
padha persis. Tanggal 29 Juni 1877
dhèwèké nulis marang Dedekind
``\emph{je le vois, mais je ne le crois pas}'';
tegesé, ``Aku weruh, nanging aku
ora percaya''. Ing ``riwayat kang
lumrah ditampa'' ngenani matématika,
ukara iki kerep dianggep nuduhaké
sepira \emph{pancen angel dipercaya}
asil-asil anyar iku kanggo
matématikawan wektu iku.
(Layang-layang kasebut disuguhaké
ing \citet{Gouvea2011}, lan bakal
kita rembug manèh ing
\olref[his][set][mythology]{sec}.
Garis gedhé buktiné Cantor ana
ing \olref[his][set][cantorplane]{sec}.)

Kasilé Cantor menehi inspirasi marang Peano
kanggo nimbang apa garis bisa dipetakaké
kanthi \emph{kontinu} menyang sawijining
bidang. Iki bakal dadi \emph{kurva kang
ngebaki ruang}. Ing \citeyear{Peano1890},
Peano mbangun kurva kaya mangkono.
Bab iki pancèn nglawan intuisi:
Euclid ndhéfinisikaké garis minangka
``dawa tanpa ambané'' (Buku I,
Dhéfinisi 2), nanging Peano nuduhaké
yèn kanthi nggulung garis kanthi
cara kang trep, dawané bisa
malih dadi ambané. Ing
\citeyear{Hilbert1891}, Hilbert
ngandharaké kurva kang ngebaki
ruang lan luwih gampang
dibayangaké, bebarengan
karo sawatara gambar.
Kurva iku dibangun
minangka runtutan;
iki enem tataran
kapisan pambanguné:
\begin{center}
\begin{tikzpicture}
\begin{scope}[xshift=20pt, yshift=20pt]
	\draw [oldiagcolorC, l-system={Hilbert curve, step=40pt, angle=90, axiom=L, order=1}] lindenmayer system;
\end{scope}
\begin{scope}[xshift=110pt, yshift=10pt]
	\draw [oldiagcolorC, l-system={Hilbert curve, step=20pt, angle=90, axiom=L, order=2}] lindenmayer system;
\end{scope}
\begin{scope}[xshift=205pt, yshift=5pt]
	\draw [oldiagcolorC, l-system={Hilbert curve, step=10pt, angle=90, axiom=L, order=3}] lindenmayer system;
\end{scope}
\begin{scope}[xshift=2.5pt, yshift=-97.5pt]
	\draw [oldiagcolorC, l-system={Hilbert curve, step=5pt, angle=90, axiom=L, order=4}] lindenmayer system;
\end{scope}
\begin{scope}[xshift=101.25pt, yshift=-98.75pt]
	\draw [oldiagcolorC, l-system={Hilbert curve, step=2.5pt, angle=90, axiom=L, order=5}] lindenmayer system;
\end{scope}
\begin{scope}[xshift=200.625pt, yshift=-99.375pt]
	\draw [oldiagcolorC, l-system={Hilbert curve, step=1.25pt, angle=90, axiom=L, order=6}] lindenmayer system;
\end{scope}
\draw[gray] (0pt,0pt) rectangle (80pt, 80pt);
\draw[gray] (100pt,0pt) rectangle (180pt, 80pt);
\draw[gray] (200pt,0pt) rectangle (280pt, 80pt);
\draw[gray] (0pt,-100pt) rectangle (80pt, -20pt);
\draw[gray] (100pt,-100pt) rectangle (180pt, -20pt);
\draw[gray] (200pt,-100pt) rectangle (280pt, -20pt);
\end{tikzpicture}  
\end{center}
Ing limit---gagasan kang nalika iku
wis didhéfinisikaké kanthi ketat---
sakabehé kothak kebak warna
!!{colorC}. Kajaba iku,
kurva Hilbert kontinu ing
saben titik nanging ora
bisa didiferensiasi ing titik
endi waé; kanthi intuisi,
amarga ing limit tanpa wates
fungsi iku malih arah kanthi
dadakan ing saben wektu.
(Konstruksi Hilbert bakal
kita andharaké luwih rinci
ing \olref[his][set][hilbertcurve]{sec}.)

Apa waé akibaté, konstruksi
géomètri kang ``patologis'' iki
dianggep alesan kanggo mamang
marang intuisi géomètri.
Konstruksi mau meh dadi
\emph{propaganda} kanggo
cara anyar nindakaké
matématika, kang pungkasané
tumeka ing téori himpunan.
Nalika mitos bab sajarah
bidhang iki dibangun
ing tembé, kerep diulang
yèn asil-asil iki
padha-padha ketat
lan nggumunaké banget.
Mula gunané rangkep:
minangka pepéling supaya
ora gumantung marang
intuisi géomètri,
lan minangka bukti
yèn cara mikir anyar
bisa ngasilaké
akèh prakara.

\end{document}
